% Replacement central-range proof for Section 7.
% This file is inserted inside the proof of Lemma 7.1.

\displayheading{Central range}

Write
\[\begin{gathered}
p_i=k_i/k_{\mathrm{co}},\quad y_i=\ell_i/k_{\mathrm{co}},\quad
 Y=\sum_i y_i,\quad I=\sum_i i y_i,\\
z_i=p_i-y_i,\quad R=\sum_i z_i=1-Y,\quad
 I_r=\sum_i i z_i=T-I.
\end{gathered}
\tag{7.12}\]
The residual vertex fraction is
\begin{equation*}
 \rho=\frac{n-m}{n}
 =R+\frac{I-TY}{\alpha-T}.
 \tag{7.13}
\end{equation*}

\displayheading{Uniform Stirling extraction}
Before approximating any factorial, equations (7.2) and (7.3) give the exact
identity
\[
\begin{split}
\log D(\ell)={}&
2\sum_i\{\log(k_i!)-\log(\ell_i!)-\log((k_i-\ell_i)!)\}\\
&-\ell_\bullet\log 2-\log(n!)+\log((n-m)!)
 +\sum_i\ell_i\log(u_i!)+\sum_i\log(\ell_i!)
 +(\log 2)b_\ell.
\end{split}
\]
We use the uniform form of Stirling's estimate
\[
  \log(t!)=t\log t-t+\sigma(t),
  \qquad
  |\sigma(t)|\le C_{\mathrm S}\log(t+2)
  \qquad(t\in\mathbb N_0),
\]
with the convention $0\log0=0$. Only a fixed number of factorials occurs in
the preceding identity, and every argument is at most $n$. Therefore the total
remainder is $O(\!\log n)$ uniformly, including when some $\ell_i$ or
$k_i-\ell_i$ is zero. Using $n=k_{\mathrm{co}}(\alpha-T)$ gives
\begin{equation*}
\begin{split}
 \log D(\ell)={}&n\rho\log\rho\\
 &+k_{\mathrm{co}}\sum_i\{2p_i\log p_i
     -2(p_i-y_i)\log(p_i-y_i)
     -y_i\log y_i-y_i+y_iE_i\}
     +O(\!\log n),
\end{split}
 \tag{7.14}
\end{equation*}
where
\begin{equation*}
 E_i=\log k_{\mathrm{co}}+\log(u_i!)+u_i-u_i\log n
       +(\log 2)\binom{u_i}{2}-\log 2.
 \tag{7.15}
\end{equation*}

The entropy term in (7.14) is uniform even at the boundary of the coordinate
box. By (5.15) and the bounded tangent-rounding displacement, the exact integer
profile satisfies $p_i\ge c_p>0$ after decreasing $c_p$ once. Hence there is an
absolute $C_{\mathrm H}$ such that, for $c_p\le p\le1$ and $0\le y\le p$,
\begin{equation*}
  2p\log p-2(p-y)\log(p-y)-y\log y-y
  \le C_{\mathrm H}y\log\!\left(\frac{\mathrm e}{y}\right).
  \tag{7.14a}
\end{equation*}
Indeed, when $y\le p/2$, the mean-value theorem bounds
$p\log p-(p-y)\log(p-y)$ by $Cy$, while $-y\log y$ supplies the only
singular term. When $y\ge p/2$, one has $y\ge c_p/2$, and the assertion follows
by compactness. Finally,
\[
  \sum_i y_i\log\!\left(\frac{\mathrm e}{y_i}\right)
  \le
  Y\log\!\left(\frac{4\mathrm e}{Y}\right),
\]
by the entropy bound on four coordinates. This is
$O(Y\log(\mathrm e/Y))$ uniformly.

Since $u_{i+1}=u_i-1$, exact subtraction in (7.15) gives
\begin{equation*}
\begin{aligned}
 E_{i+1}-E_i
 &=\log n-\log(\alpha-i)-(\log 2)(\alpha-i)+\log 2-1\\
 &=-\frac{\log 2}{2}\alpha+O(1),
\end{aligned}
 \tag{7.16}
\end{equation*}
uniformly for $2\le i\le4$. For the second line, use
$(\log 2)\alpha/2=\log n-\log\log n+O(1)$ and
$\log(\alpha-i)=\log\log n+O(1)$.

Let $\bar E=\sum_i p_iE_i$. Applying Stirling once to the complete signed
first moment gives
\begin{equation*}
 -\frac1{k_{\mathrm{co}}}
   \log \Exp{Z_{\mathbf k}^{\mathrm{sgn}}}
 =\sum_i p_i\log p_i-1+\bar E+o(1).
 \tag{7.17}
\end{equation*}
The phase-uniform margin (5.19) implies $\bar E\le C_E$. To expose the affine
structure in (7.16), set
\[
  F_i:=E_i+\frac{(\log 2)\alpha}{2}i.
\]
Then $F_{i+1}-F_i=O(1)$. Because the support has only four consecutive
indices,
\[
  F_i=\sum_jp_jF_j+O(1)
     =\bar E+\frac{(\log 2)\alpha}{2}T+O(1).
\]
Consequently
\[
  E_i=\bar E-\frac{(\log 2)\alpha}{2}(i-T)+O(1),
\]
and therefore
\begin{equation*}
 \sum_i y_iE_i
 \le\frac{(\log 2)\alpha}{2}(TY-I)+O(Y).
 \tag{7.18}
\end{equation*}
Only the upper bound on $\bar E$ is used here.

It remains to replace the true residual fraction $\rho$ by the profile
fraction $R$. Since $2\le i,T\le5$,
\[
  |I-TY|=\left|\sum_i(i-T)y_i\right|\le3Y,
\]
so (7.13) gives
\[
  |\rho-R|\le\frac{3Y}{\alpha-T}=O(Y/\alpha).
\]
If $\rho\ge1/32$, then $Y\le1$ and the preceding bound implies
$R\ge1/64$ for all sufficiently large $n$. Hence $x\mapsto x\log x$ is
uniformly Lipschitz on the interval between $R$ and $\rho$. Moreover,
$-R\log R\le1-R=Y$. Thus
\[
\begin{aligned}
&\left|n\rho\log\rho
      -k_{\mathrm{co}}\alpha R\log R\right|\\
&\qquad\le
 k_{\mathrm{co}}(\alpha-T)C|\rho-R|
 +k_{\mathrm{co}}T|R\log R|
 \le Ck_{\mathrm{co}}Y,
\end{aligned}
\]
which proves
\begin{equation*}
 n\rho\log\rho
 =k_{\mathrm{co}}\alpha R\log R+O(k_{\mathrm{co}}Y).
 \tag{7.19}
\end{equation*}
Combining (7.14), (7.14a), (7.18), and (7.19) yields one absolute constant
$C$ such that
\begin{equation*}
 \log D(\ell)
 \le k_{\mathrm{co}}\alpha\Phi_T(z)
     +Ck_{\mathrm{co}}Y\log\!\left(\frac{\mathrm e}{Y}\right)
     +C\log n,
 \tag{7.20}
\end{equation*}
where
\begin{equation*}
 \Phi_T(z)=R\log R+\frac{\log 2}{2}(I_r-TR).
 \tag{7.21}
\end{equation*}
All constants and eventuality thresholds in this extraction are independent of
the phase.

\displayheading{Uniform rate negativity}
The four-deficit geometry gives
\begin{equation*}
 I_r-TR\le(5-T)R,
 \qquad
 I_r-TR=\sum_i(T-i)y_i\le(T-2)(1-R).
 \tag{7.22}
\end{equation*}
Put $q=\log 2$ and recall that $Y=1-R$. Multiplying the first inequality in
(7.22) by $Y$, the second by $R$, and adding gives
\begin{equation*}
 I_r-TR\le3RY.
 \tag{7.22a}
 \label{eq:partial-diagonal-combined-structural-v3}
\end{equation*}
For $0<R\le1$,
\[
 \log R\le\frac{2(R-1)}{R+1}.
\]
Indeed, for $x=(1-R)/R\ge0$, the function
$\log(1+x)-2x/(2+x)$ has derivative
$x^2/((1+x)(2+x)^2)\ge0$ and vanishes at zero. The positive expansion for
$\log 2$ gives
\[
  \frac23<q<\frac7{10}.
\]

If $1/64\le R\le3/4$, then
\[
\begin{aligned}
 \Phi_T
 &\le-\frac{2R}{1+R}Y+\frac{21}{20}RY\\
 &=-\left(\frac{2}{1+R}-\frac{21}{20}\right)RY
 \le-\frac{13}{8960}Y.
\end{aligned}
\]
Here $2/(1+R)\ge8/7$ and $R\ge1/64$. If $3/4\le R\le1$, then the phase
corridor gives $T<4$, so the second inequality in (7.22) implies
$I_r-TR\le2Y$. Since $\log R\le R-1=-Y$,
\[
  \Phi_T\le-RY+qY=(q-R)Y\le-\frac1{20}Y.
\]
Thus, throughout the scalar range,
\begin{equation*}
  \Phi_T(z)\le-\frac{Y}{5000}.
  \tag{7.23}
\end{equation*}
The rational inequalities in this two-range ledger are independently checked
by \path{check_partial_diagonal_rate_v3.py}; the logarithmic inequality and the
four-deficit reduction are proved above.

\displayheading{Uniform summation over the central range}
Assume
\begin{equation*}
  m>\eta n,
  \qquad
  n-m>n/32.
  \tag{7.24}
\end{equation*}
Then $\rho>1/32$, and
\[
  \frac mn=\frac{\alpha Y-I}{\alpha-T}
  \le\frac{\alpha Y}{\alpha-T}
\]
implies
\[
  Y>\eta\frac{\alpha-T}{\alpha}\ge\frac\eta2
\]
for all sufficiently large $n$. Define the deterministic error ratio
\begin{equation*}
  \varepsilon_n^{\mathrm{diag}}
  :=
  \frac{C\log(2\mathrm e/\eta)}{\alpha}
  +
  \frac{2C\log n}{k_{\mathrm{co}}\alpha\eta}.
  \tag{7.24a}
\end{equation*}
Since
\[
  \alpha=\Theta(\log n),
  \qquad
  k_{\mathrm{co}}=\Theta(n/\log n),
  \qquad
  \eta=\Theta(\log\log n/\log n),
\]
one has $\varepsilon_n^{\mathrm{diag}}\to0$, uniformly in the phase. Because
$Y\ge\eta/2$, equations (7.20) and (7.23) give
\[
  \log D(\ell)
  \le
  -k_{\mathrm{co}}\alpha Y
   \left(\frac1{5000}-\varepsilon_n^{\mathrm{diag}}\right).
\]
For all sufficiently large $n$, the parenthesis is at least $1/10000$.
Moreover, $\alpha\eta=\Theta(\log\log n)$. Hence there is a phase-independent
$c>0$ such that
\begin{equation*}
  D(\ell)
  \le
  \exp\{-c k_{\mathrm{co}}\log\log n\}.
  \tag{7.25}
\end{equation*}
There are at most $(k_{\mathrm{co}}+1)^4=\exp(O(\log n))$ subprofiles, whereas
$k_{\mathrm{co}}\log\log n\gg\log n$. Their total central-range contribution
is therefore $o(1)$, with one eventuality threshold for the complete phase.
